\chapter{Sprint 2 --- Day 2 (partial): Performance optimizations}

\textbf{Cập nhật 2026-04-30:} Method \texttt{\_get\_accessible\_branch\_partner\_ids} đã đổi tên thành \texttt{\_get\_accessible\_franchise\_ids} (trả tuple \texttt{wujia.franchise.management.id} thay cho \texttt{res.partner.id}). Cache invalidation hooks vẫn ở CRUD của member model (chỉ đổi tên model \texttt{wujia.branch.member} → \texttt{wujia.franchise.member}). Logic + chiến thuật cache không thay đổi.

\section{Phạm vi Day 2 (đã làm)}

Theo BA chỉ đạo "tập trung tối ưu trước, các phần khác từ từ", Day 2 thu hẹp scope còn:
\begin{itemize}
    \item Cài \texttt{queue\_job} (OCA branch 19.0) làm tool async cho tương lai
    \item Fix \textbf{P0-1}: Apply \texttt{@ormcache} cho \texttt{res.users.\_get\_accessible\_franchise\_ids} (rename từ \texttt{\_branch\_partner\_ids})
    \item Fix \textbf{P0-2}: Store + index \texttt{is\_currently\_valid} trên \texttt{wujia.franchise.member} (rename từ \texttt{wujia.branch.member})
\end{itemize}

Defer Day 2 còn lại (service layer, bus emit cho sale.order, unit tests) sang lần làm sau.

\section{Setup \texttt{queue\_job} (chưa dùng, sẵn sàng cho sau)}

\begin{lstlisting}[style=shell]
cd /home/huyban/odoo-dev/WujiaTea/themes
git clone --branch 19.0 --single-branch --depth 1 \
  https://github.com/OCA/queue.git oca_queue
\end{lstlisting}

Cập nhật \texttt{config/odoo.conf} \texttt{addons\_path} thêm \texttt{themes/oca\_queue}.
Install: \texttt{python odoo-bin -i queue\_job --stop-after-init}.

\section{Fix P0-1: \texttt{ormcache} cho lookup chi nhánh user}

\textbf{Vấn đề trước fix:} Mỗi câu \texttt{SELECT * FROM sale\_order} từ portal user
trigger eval record rule → gọi \texttt{user.\_get\_accessible\_branch\_partners().ids}
→ 2 query phụ (membership + partner mapped). Với 1000 user concurrent =
2000+ query/giây chỉ để check security.

\textbf{Sau fix:} Method \texttt{\_get\_accessible\_branch\_partner\_ids()} cache
per-worker. Subsequent calls trả về tuple từ memory, không hit DB.

\subsection{Code thay đổi}

File: \texttt{custom/wujia\_franchise/models/res\_users.py}

\begin{lstlisting}[style=python]
from odoo.tools import ormcache

@ormcache('self.id')
def _get_accessible_branch_partner_ids(self):
    self.ensure_one()
    return tuple(self._get_active_branch_memberships()
                 .mapped('branch_partner_id').ids)

def _get_accessible_branch_partners(self):
    self.ensure_one()
    return self.env['res.partner'].browse(
        self._get_accessible_branch_partner_ids()
    )
\end{lstlisting}

\subsection{Cache invalidation}

File: \texttt{custom/wujia\_franchise/models/wujia\_branch\_member.py}

Hook vào \texttt{create/write/unlink}:

\begin{lstlisting}[style=python]
def _invalidate_user_cache(self):
    if self:
        self.env.registry.clear_cache()  # cross-worker safe

@api.model_create_multi
def create(self, vals_list):
    records = super().create(vals_list)
    records._invalidate_user_cache()
    return records

def write(self, vals):
    res = super().write(vals)
    self._invalidate_user_cache()
    return res

def unlink(self):
    had_records = bool(self)
    res = super().unlink()
    if had_records:
        self.env.registry.clear_cache()
    return res
\end{lstlisting}

\textbf{Lưu ý:} Dùng \texttt{env.registry.clear\_cache()} thay vì
\texttt{method.clear\_cache(records)} — Odoo 19 tự signal qua DB để
invalidate cache trên mọi worker khác. Heavy hơn 1 chút nhưng cross-worker safe.

\subsection{Update record rule dùng method mới}

File: \texttt{custom/wujia\_sale/security/wujia\_sale\_rules.xml}

Đổi \texttt{user.\_get\_accessible\_branch\_partners().ids} →
\texttt{user.\_get\_accessible\_branch\_partner\_ids()} (tuple cached, không
trigger ORM search mỗi query).

\section{Fix P0-2: Store + index \texttt{is\_currently\_valid}}

\textbf{Vấn đề trước fix:} \texttt{is\_currently\_valid} là compute non-stored với
search method tự viết. Mọi search \texttt{[('is\_currently\_valid', '=', True)]}
phải iterate trong Python — không dùng được index, không scale 9000 membership.

\textbf{Sau fix:} Field stored + B-tree index → SQL search direct, query
plan dùng index scan thay vì sequential scan.

\subsection{Code thay đổi}

File: \texttt{custom/wujia\_franchise/models/wujia\_branch\_member.py}

\begin{lstlisting}[style=python]
is_currently_valid = fields.Boolean(
    compute='_compute_is_currently_valid',
    store=True,           # Đã thêm
    index=True,           # Đã thêm
)
# Bỏ method _search_is_currently_valid (không cần khi store)
\end{lstlisting}

\subsection{Cron daily recompute}

Compute phụ thuộc \texttt{today} → không tự re-trigger qua \texttt{@api.depends}.
Cần cron daily update các record vừa enter/exit hiệu lực.

File: \texttt{custom/wujia\_franchise/models/wujia\_branch\_member.py}

\begin{lstlisting}[style=python]
@api.model
def _cron_recompute_validity(self):
    from datetime import timedelta
    today = fields.Date.context_today(self)
    affected = self.search([
        '|',
        ('date_from', '=', today),               # vừa active
        ('date_to', '=', today - timedelta(days=1)),  # vừa hết hạn
    ])
    if affected:
        affected._compute_is_currently_valid()
        affected.flush_recordset()
        affected._invalidate_user_cache()
\end{lstlisting}

File mới: \texttt{custom/wujia\_franchise/data/ir\_cron\_data.xml} — cron
chạy hàng ngày interval=1 day.

\subsection{Migration}

Khi upgrade module, Odoo tự re-init compute cho mọi record → store column
populated tự động. Verify 7/7 record demo có \texttt{is\_currently\_valid=True}.

\section{Verify đã pass}

\begin{itemize}
    \item Module install/upgrade clean (0 error)
    \item Cột \texttt{is\_currently\_valid} stored, có B-tree index
    \texttt{wujia\_branch\_member\_\_is\_currently\_valid\_index}
    \item Cron \texttt{Wujia: Recompute branch member is\_currently\_valid}
    active, daily
    \item Test cache: 1st call → DB hit, 2nd call → memory hit
    \item Test invalidate: \texttt{member.active=False} → cache cleared,
    next call trả về tuple rỗng đúng
    \item Test record rule: \texttt{anh.owner} (HN-01 only) search SO chỉ
    thấy 2 đơn HN-01, không thấy đơn HCM-01 admin tạo
\end{itemize}

\section{Đo lường impact (theoretical)}

\begin{longtable}{p{4cm} p{4cm} p{6cm}}
\toprule
\textbf{Metric} & \textbf{Trước fix} & \textbf{Sau fix} \\
\midrule
\endhead
Query/page-load (record rule) & 2--3 query/user/page & 0 (cache hit ratio $>$ 90\%) \\
Search \texttt{is\_currently\_valid=True} & Sequential scan (Python iter) & Index scan (B-tree) \\
Multi-user concurrent (1000 user × 5 page-load/min) & 10k--15k query/min & 1k--2k query/min \\
Worst case @ 9k membership records & 9k row Python loop & 1 SQL query (millis) \\
\bottomrule
\end{longtable}

\section{Files đã chạm Day 2}

\begin{itemize}
    \item Mới: \texttt{themes/oca\_queue/} (clone)
    \item Mới: \texttt{custom/wujia\_franchise/data/ir\_cron\_data.xml}
    \item Edit: \texttt{config/odoo.conf} — \texttt{addons\_path} thêm \texttt{themes/oca\_queue}
    \item Edit: \texttt{custom/wujia\_franchise/\_\_manifest\_\_.py} — data thêm cron
    \item Edit: \texttt{custom/wujia\_franchise/models/res\_users.py} — \texttt{ormcache}
    \item Edit: \texttt{custom/wujia\_franchise/models/wujia\_branch\_member.py} — \texttt{store=True, index=True}; bỏ \texttt{\_search\_*}; thêm CRUD hooks invalidate cache; thêm \texttt{\_cron\_recompute\_validity}
    \item Edit: \texttt{custom/wujia\_sale/security/wujia\_sale\_rules.xml} — đổi domain dùng method tuple cached
\end{itemize}

% ===========================================================================
% ===========================================================================
